% Generated draft: requires mathematical and source review.
\documentclass[11pt]{article}
\usepackage[margin=25mm]{geometry}
\usepackage{fontspec}
\setmainfont{Latin Modern Roman}
\usepackage{amsmath,amssymb,mathtools,unicode-math}
\setmathfont{Latin Modern Math}
\usepackage{enumitem,needspace,xurl,hyperref}
\setlength{\parindent}{0pt}
\setlength{\parskip}{5pt}
\setlength{\emergencystretch}{3em}
\newcommand{\N}{\mathbb{N}}
\newcommand{\Z}{\mathbb{Z}}
\newcommand{\Q}{\mathbb{Q}}
\newcommand{\R}{\mathbb{R}}
\newcommand{\C}{\mathbb{C}}
\newcommand{\F}{\mathbb{F}}
\newcommand{\A}{\mathbb{A}}
\newcommand{\PP}{\mathbb{P}}
\newcommand{\E}{\mathbb{E}}
\newcommand{\eps}{\varepsilon}
\newcommand{\dd}{\,\mathrm{d}}
\begin{document}
\section*{UnsolvedMath ID: 800678}
\textbf{Source problem code:} catalogue entry
\subsection{Definitions and mathematical statement}
Let \(\PP^2(\C)\) denote the complex projective plane, and let \(\mathscr F\) be a singular holomorphic foliation by complex curves on \(\PP^2(\C)\). Write
\[
S=\operatorname{Sing}(\mathscr F)
\]
for its singular set, which is assumed to be finite. On \(\PP^2(\C)\setminus S\), the foliation determines a holomorphic tangent distribution of complex rank one. A leaf is a maximal connected immersed complex curve tangent to this distribution. For a leaf \(L\), let \(\overline L\) denote its closure in the usual topology of \(\PP^2(\C)\).

The question is whether, for every such foliation \(\mathscr F\) and every leaf \(L\) of \(\mathscr F\),
\[
\overline L\cap\operatorname{Sing}(\mathscr F)\ne\varnothing.
\]
Here the complex ground field, the complex dimension of the leaves, and the usual topology for closure make explicit the conventional interpretation of the supplied holomorphic-foliation statement.
\subsection{Short English statement}
Must the closure of every leaf of every holomorphic foliation on the complex projective plane with finite singular set contain a singular point?
\subsection{Sources}
\begin{itemize}
\item Local UnsolvedMath registry, permanent ID \(800678\), source problem code \texttt{catalogue entry}, titled ``Holomorphic Foliation Theory --- Problem 5.''
\item Dinh and Sibony, \emph{Some Open Problems on Holomorphic Foliation Theory} (2018), Problem 5. Supplied source URL: \url{https://arxiv.org/abs/1807.05640}. The registry attributes the problem to Dinh and Sibony, with proposed year 2018.
\item Supplied provenance records access on 2026-10-16 and extraction by \texttt{arxiv\_latex}. These are supplied provenance claims; the source has not been independently visited or verified for this draft.
\item The supplied registry labels the problem ``open.'' This is a supplied status label, not an independently verified assessment of the literature.
\end{itemize}
\end{document}
